% Section 5.2, from lectures/L05/README.md: the objectives, the evaluation questions, and the next
% lecture.

\section{Review}
\label{c5:sec:review}

\needspace{6\baselineskip}
\subsection*{What you should be able to do now}
After this chapter, you should be able to:
\begin{itemize}
  \item \textbf{Explain why a single-cycle error pulse cannot be polled}, and implement
    pulse-to-level latching with a defined set and clear pair for each \code{ERROR_FLAGS} bit, while
    deriving every \code{STATUS} bit combinationally from the FIFO flags instead.
  \item \textbf{Explain why popping a FIFO on read would need write-like commit-and-abort
    discipline}, and why splitting the pure read (\code{RX_DATA}) from the explicit pop
    (\code{RX_POP}) avoids it.
  \item \textbf{Instantiate the register bank into \code{uart_top} positionally}, completing the
    peripheral so that it passes the system testbench.
\end{itemize}

\needspace{6\baselineskip}
\subsection*{Questions to test yourself}
You should be able to explain:
\begin{enumerate}
  \item Why can \code{STATUS} bit 1 (RX valid) not simply be wired to the receiver's ``byte ready''
    pulse?
  \item A \code{TX_DATA} write arrives while \code{STATUS} bit 0 (TX ready) is clear: what happens,
    and why is dropping the byte the right choice here?
  \item \code{SS} rises after the third byte of an \code{RX_POP} write: what must the peripheral do,
    and which part of the transport contract guarantees it?
  \item If reading \code{RX_DATA} returns a byte but never advances the FIFO, what did the driver
    forget, and what would go wrong if \code{RX_DATA} popped on its own?
\end{enumerate}

\needspace{6\baselineskip}
\subsection*{Looking ahead}
The FPGA half is done. The driver stack begins in \chapref{6}, on the host, with no hardware: the
\code{driver::uart::Interface} and its \code{driver::transport::Interface} seam, then a first
\code{driver::uart::Stub} written against that interface.
